\documentclass[12pt,]{article}
\usepackage{lmodern}
\usepackage{amssymb,amsmath}
\usepackage{ifxetex,ifluatex}
\usepackage{fixltx2e} % provides \textsubscript
\ifnum 0\ifxetex 1\fi\ifluatex 1\fi=0 % if pdftex
  \usepackage[T1]{fontenc}
  \usepackage[utf8]{inputenc}
\else % if luatex or xelatex
  \ifxetex
    \usepackage{mathspec}
  \else
    \usepackage{fontspec}
  \fi
  \defaultfontfeatures{Ligatures=TeX,Scale=MatchLowercase}
\fi
% use upquote if available, for straight quotes in verbatim environments
\IfFileExists{upquote.sty}{\usepackage{upquote}}{}
% use microtype if available
\IfFileExists{microtype.sty}{%
\usepackage{microtype}
\UseMicrotypeSet[protrusion]{basicmath} % disable protrusion for tt fonts
}{}
\usepackage[margin=1in]{geometry}
\usepackage{hyperref}
\hypersetup{unicode=true,
            pdftitle={Jacobi Eigen in R and C with Lower Triangular Column-wise Compact Storage},
            pdfauthor={Jan de Leeuw},
            pdfborder={0 0 0},
            breaklinks=true}
\urlstyle{same}  % don't use monospace font for urls
\usepackage{color}
\usepackage{fancyvrb}
\newcommand{\VerbBar}{|}
\newcommand{\VERB}{\Verb[commandchars=\\\{\}]}
\DefineVerbatimEnvironment{Highlighting}{Verbatim}{commandchars=\\\{\}}
% Add ',fontsize=\small' for more characters per line
\usepackage{framed}
\definecolor{shadecolor}{RGB}{248,248,248}
\newenvironment{Shaded}{\begin{snugshade}}{\end{snugshade}}
\newcommand{\KeywordTok}[1]{\textcolor[rgb]{0.13,0.29,0.53}{\textbf{#1}}}
\newcommand{\DataTypeTok}[1]{\textcolor[rgb]{0.13,0.29,0.53}{#1}}
\newcommand{\DecValTok}[1]{\textcolor[rgb]{0.00,0.00,0.81}{#1}}
\newcommand{\BaseNTok}[1]{\textcolor[rgb]{0.00,0.00,0.81}{#1}}
\newcommand{\FloatTok}[1]{\textcolor[rgb]{0.00,0.00,0.81}{#1}}
\newcommand{\ConstantTok}[1]{\textcolor[rgb]{0.00,0.00,0.00}{#1}}
\newcommand{\CharTok}[1]{\textcolor[rgb]{0.31,0.60,0.02}{#1}}
\newcommand{\SpecialCharTok}[1]{\textcolor[rgb]{0.00,0.00,0.00}{#1}}
\newcommand{\StringTok}[1]{\textcolor[rgb]{0.31,0.60,0.02}{#1}}
\newcommand{\VerbatimStringTok}[1]{\textcolor[rgb]{0.31,0.60,0.02}{#1}}
\newcommand{\SpecialStringTok}[1]{\textcolor[rgb]{0.31,0.60,0.02}{#1}}
\newcommand{\ImportTok}[1]{#1}
\newcommand{\CommentTok}[1]{\textcolor[rgb]{0.56,0.35,0.01}{\textit{#1}}}
\newcommand{\DocumentationTok}[1]{\textcolor[rgb]{0.56,0.35,0.01}{\textbf{\textit{#1}}}}
\newcommand{\AnnotationTok}[1]{\textcolor[rgb]{0.56,0.35,0.01}{\textbf{\textit{#1}}}}
\newcommand{\CommentVarTok}[1]{\textcolor[rgb]{0.56,0.35,0.01}{\textbf{\textit{#1}}}}
\newcommand{\OtherTok}[1]{\textcolor[rgb]{0.56,0.35,0.01}{#1}}
\newcommand{\FunctionTok}[1]{\textcolor[rgb]{0.00,0.00,0.00}{#1}}
\newcommand{\VariableTok}[1]{\textcolor[rgb]{0.00,0.00,0.00}{#1}}
\newcommand{\ControlFlowTok}[1]{\textcolor[rgb]{0.13,0.29,0.53}{\textbf{#1}}}
\newcommand{\OperatorTok}[1]{\textcolor[rgb]{0.81,0.36,0.00}{\textbf{#1}}}
\newcommand{\BuiltInTok}[1]{#1}
\newcommand{\ExtensionTok}[1]{#1}
\newcommand{\PreprocessorTok}[1]{\textcolor[rgb]{0.56,0.35,0.01}{\textit{#1}}}
\newcommand{\AttributeTok}[1]{\textcolor[rgb]{0.77,0.63,0.00}{#1}}
\newcommand{\RegionMarkerTok}[1]{#1}
\newcommand{\InformationTok}[1]{\textcolor[rgb]{0.56,0.35,0.01}{\textbf{\textit{#1}}}}
\newcommand{\WarningTok}[1]{\textcolor[rgb]{0.56,0.35,0.01}{\textbf{\textit{#1}}}}
\newcommand{\AlertTok}[1]{\textcolor[rgb]{0.94,0.16,0.16}{#1}}
\newcommand{\ErrorTok}[1]{\textcolor[rgb]{0.64,0.00,0.00}{\textbf{#1}}}
\newcommand{\NormalTok}[1]{#1}
\usepackage{graphicx,grffile}
\makeatletter
\def\maxwidth{\ifdim\Gin@nat@width>\linewidth\linewidth\else\Gin@nat@width\fi}
\def\maxheight{\ifdim\Gin@nat@height>\textheight\textheight\else\Gin@nat@height\fi}
\makeatother
% Scale images if necessary, so that they will not overflow the page
% margins by default, and it is still possible to overwrite the defaults
% using explicit options in \includegraphics[width, height, ...]{}
\setkeys{Gin}{width=\maxwidth,height=\maxheight,keepaspectratio}
\IfFileExists{parskip.sty}{%
\usepackage{parskip}
}{% else
\setlength{\parindent}{0pt}
\setlength{\parskip}{6pt plus 2pt minus 1pt}
}
\setlength{\emergencystretch}{3em}  % prevent overfull lines
\providecommand{\tightlist}{%
  \setlength{\itemsep}{0pt}\setlength{\parskip}{0pt}}
\setcounter{secnumdepth}{5}
% Redefines (sub)paragraphs to behave more like sections
\ifx\paragraph\undefined\else
\let\oldparagraph\paragraph
\renewcommand{\paragraph}[1]{\oldparagraph{#1}\mbox{}}
\fi
\ifx\subparagraph\undefined\else
\let\oldsubparagraph\subparagraph
\renewcommand{\subparagraph}[1]{\oldsubparagraph{#1}\mbox{}}
\fi

%%% Use protect on footnotes to avoid problems with footnotes in titles
\let\rmarkdownfootnote\footnote%
\def\footnote{\protect\rmarkdownfootnote}

%%% Change title format to be more compact
\usepackage{titling}

% Create subtitle command for use in maketitle
\providecommand{\subtitle}[1]{
  \posttitle{
    \begin{center}\large#1\end{center}
    }
}

\setlength{\droptitle}{-2em}

  \title{Jacobi Eigen in R and C with Lower Triangular Column-wise Compact
Storage}
    \pretitle{\vspace{\droptitle}\centering\huge}
  \posttitle{\par}
    \author{Jan de Leeuw}
    \preauthor{\centering\large\emph}
  \postauthor{\par}
      \predate{\centering\large\emph}
  \postdate{\par}
    \date{Version 01, July 11, 2017}


\begin{document}
\maketitle
\begin{abstract}
The Jacobi method for computing eigenvalues and eigenvectors of a
symmetric matrix is implemented in C using column-wise compact storage
of the lower triangle. The complied C code can be loaded into R using
the .C() interface. We compare the C implementation with an earlier
version in pure R, and with the built-in eigen function in R.
\end{abstract}

{
\setcounter{tocdepth}{3}
\tableofcontents
}
Note: This is a working paper which will be expanded/updated frequently.
All suggestions for improvement are welcome. The directory
\href{http://deleeuwpdx.net/pubfolders/jacobi}{deleeuwpdx.net/pubfolders/jacobi}
has a pdf version, the complete Rmd file with all code chunks, the R and
C code, and the bib file.

\section{Introduction}\label{introduction}

In multivariate analysis we often deal with symmetric (and positive
semi-definite) matrices. Computer code, both in R or in C, often does
not take symmetry into account when storing such objects. Thus a
symmetric matrix is stored in memory like any other matrix, with the
obvious redundancy all all off-diagonal elements are stired twice. This
may be smart from a computational and programming point of view, because
all operations on matrices are available without change. But I, for one,
have always had a gnawing feeling of unease, because of this blemish of
redundancy.

There are, of course, ways to store a symmetric matrix more efficiently.
We can store either the upper or lower triangle, and we can store either
one row-wise or column-wise. In this paper we have chosen to use the
column-wise lower triangle, and we apply the classical iterative Jacobi
method to compute eigenvalues and eignevectors to this stored triangle.

\section{Implementation}\label{implementation}

The basic computations are in C. This is done mostly for speed, but also
because I like to think of R as a wrapper for compiled C code. The fact
that our C code is going to be called from R makes it natural to store
our matrices column-wise.

The C code has a number of idiosyncracies. I use small static inline
functions, defined and declared in the header file, to compute locations
in memory (i.e.~to map row and column matrix indices to locations in the
triangle, stored column-wise as vector). The C code is written in such a
way that it can easily be used on systems where R is not available.
There are no R includes and no R functions called. But on the other hand
the code follows most of the conventions of the .C() interface in R --
i.e.~all functions return void, and all arguments are passed as
pointers. There is also no I/O and dynamic memory allocation, because I
generally prefer to delegate that to the calling system, i.e.~our case
to R.

The implementation of the Jacobi method does not have any bells and
whistles. It cycles through all off-diagonal elements in order. It does
not incorporate some well-known ways to use parallelism and sparseness.
We have closely followed an earlier implementation in pure R (De Leeuw
and Ferrari (2008)). But all in all, I needed something like this for a
new R/C implement of smacof which uses the same approach to storing
symmetric matrices. There is no attempt to compete with Lapack in terms
of speed or robustness. Of course jacobi() has a parameter which
determines precision, and this may be useful if only an approximate
diagonalization is needed.

\section{Example}\label{example}

We first try eigen(), jacobi(), and the R implementation jacobiR() from
({\textbf{???}}) to verify the output is the same. All three routines
are applied to a Hilbert matrix of order 4.

\begin{Shaded}
\begin{Highlighting}[]
\NormalTok{h <-}\StringTok{ }\DecValTok{1}\OperatorTok{/}\NormalTok{(}\KeywordTok{outer}\NormalTok{(}\DecValTok{1}\OperatorTok{:}\DecValTok{4}\NormalTok{,}\DecValTok{1}\OperatorTok{:}\DecValTok{4}\NormalTok{,}\StringTok{"+"}\NormalTok{) }\OperatorTok{-}\StringTok{ }\DecValTok{1}\NormalTok{)}
\NormalTok{g <-}\StringTok{ }\KeywordTok{matrix2triangle}\NormalTok{ (h)}
\KeywordTok{eigen}\NormalTok{(h)}
\end{Highlighting}
\end{Shaded}

\begin{verbatim}
## eigen() decomposition
## $values
## [1] 1.500214280e+00 1.691412202e-01 6.738273606e-03 9.670230402e-05
## 
## $vectors
##              [,1]          [,2]          [,3]           [,4]
## [1,] 0.7926082912  0.5820756995 -0.1791862905 -0.02919332316
## [2,] 0.4519231209 -0.3705021851  0.7419177906  0.32871205576
## [3,] 0.3224163986 -0.5095786345 -0.1002281369 -0.79141114583
## [4,] 0.2521611697 -0.5140482722 -0.6382825282  0.51455275000
\end{verbatim}

\begin{Shaded}
\begin{Highlighting}[]
\KeywordTok{jacobi}\NormalTok{(g)}
\end{Highlighting}
\end{Shaded}

\begin{verbatim}
## $eval
## [1] 1.500214280e+00 1.691412202e-01 6.738273606e-03 9.670230402e-05
## 
## $evec
##              [,1]          [,2]          [,3]           [,4]
## [1,] 0.7926082911 -0.5820756995  0.1791862906 -0.02919332318
## [2,] 0.4519231210  0.3705021851 -0.7419177906  0.32871205578
## [3,] 0.3224163986  0.5095786345  0.1002281370 -0.79141114581
## [4,] 0.2521611696  0.5140482722  0.6382825282  0.51455275003
\end{verbatim}

\begin{Shaded}
\begin{Highlighting}[]
\KeywordTok{jacobiR}\NormalTok{(h)}
\end{Highlighting}
\end{Shaded}

\begin{verbatim}
## $values
## [1] 1.500214280e+00 1.691412202e-01 6.738273606e-03 9.670230402e-05
## 
## $vectors
##              [,1]          [,2]          [,3]           [,4]
## [1,] 0.7926082913  0.5820756996  0.1791862916  0.02919331096
## [2,] 0.4519231201 -0.3705021850 -0.7419178134 -0.32871200562
## [3,] 0.3224163988 -0.5095786345  0.1002281901  0.79141113902
## [4,] 0.2521611704 -0.5140482722  0.6382824931 -0.51455279320
\end{verbatim}

We use a Hilbert matrix of order 100 to compare the running speed of
eigen(), eigen() with symmetric=TRUE, jacobi(), and jacobiR(). For this
example, and for the chosen level of precision, eigen() without
symmetric is takes about 1.5 times the time that eigen() with symmetric
does, jacobi() takes four to five times as long as eigen(), and
jacobiR() takes at least 100 times the time of jacobiC().

\begin{Shaded}
\begin{Highlighting}[]
\NormalTok{h <-}\StringTok{ }\DecValTok{1}\OperatorTok{/}\NormalTok{(}\KeywordTok{outer}\NormalTok{(}\DecValTok{1}\OperatorTok{:}\DecValTok{100}\NormalTok{,}\DecValTok{1}\OperatorTok{:}\DecValTok{100}\NormalTok{,}\StringTok{"+"}\NormalTok{) }\OperatorTok{-}\StringTok{ }\DecValTok{1}\NormalTok{)}
\NormalTok{g <-}\StringTok{ }\KeywordTok{matrix2triangle}\NormalTok{ (h)}
\KeywordTok{microbenchmark}\NormalTok{ (}\KeywordTok{eigen}\NormalTok{(h), }\KeywordTok{eigen}\NormalTok{(h, }\DataTypeTok{symmetric =} \OtherTok{TRUE}\NormalTok{), }\KeywordTok{jacobi}\NormalTok{(g), }\KeywordTok{jacobiR}\NormalTok{(h))}
\end{Highlighting}
\end{Shaded}

\begin{verbatim}
## Unit: microseconds
##                        expr        min          lq         mean
##                    eigen(h)    997.528   1096.7200   1558.17462
##  eigen(h, symmetric = TRUE)    680.777    713.2980    914.11742
##                   jacobi(g)   3459.693   3615.8415   3834.97385
##                  jacobiR(h) 428195.412 447114.0620 464897.94215
##       median          uq        max neval cld
##    1188.6970   1331.7410   5108.599   100  a 
##     752.0980    870.7415   4945.969   100  a 
##    3708.8975   3885.5455   6871.676   100  a 
##  459867.4815 482148.1105 515549.009   100   b
\end{verbatim}

\section{Appendix: Code}\label{appendix-code}

\subsection{R code}\label{r-code}

\begin{Shaded}
\begin{Highlighting}[]
\KeywordTok{dyn.load}\NormalTok{(}\StringTok{"jacobi.so"}\NormalTok{)}

\NormalTok{jacobi <-}\StringTok{ }\ControlFlowTok{function}\NormalTok{ (a,}
                    \DataTypeTok{itmax =} \DecValTok{10}\NormalTok{,}
                    \DataTypeTok{eps =} \FloatTok{1e-6}\NormalTok{,}
                    \DataTypeTok{verbose =} \OtherTok{FALSE}\NormalTok{) \{}
\NormalTok{  m <-}\StringTok{ }\KeywordTok{length}\NormalTok{ (a)}
\NormalTok{  n <-}\StringTok{ }\KeywordTok{round}\NormalTok{ ((}\KeywordTok{sqrt}\NormalTok{ (}\DecValTok{1} \OperatorTok{+}\StringTok{ }\DecValTok{8} \OperatorTok{*}\StringTok{ }\NormalTok{m) }\OperatorTok{-}\StringTok{ }\DecValTok{1}\NormalTok{) }\OperatorTok{/}\StringTok{ }\DecValTok{2}\NormalTok{)}
\NormalTok{  r <-}\StringTok{ }\OperatorTok{-}\NormalTok{((}\DecValTok{1}\OperatorTok{:}\NormalTok{n) }\OperatorTok{^}\StringTok{ }\DecValTok{2}\NormalTok{) }\OperatorTok{/}\StringTok{ }\DecValTok{2} \OperatorTok{+}\StringTok{ }\NormalTok{(}\DecValTok{2} \OperatorTok{*}\StringTok{ }\NormalTok{n }\OperatorTok{+}\StringTok{ }\DecValTok{3}\NormalTok{) }\OperatorTok{*}\StringTok{ }\NormalTok{(}\DecValTok{1}\OperatorTok{:}\NormalTok{n) }\OperatorTok{/}\StringTok{ }\DecValTok{2} \OperatorTok{-}\StringTok{ }\NormalTok{n}
\NormalTok{  h <-}
\StringTok{    }\KeywordTok{.C}\NormalTok{(}
      \StringTok{"jacobiC"}\NormalTok{,}
      \KeywordTok{as.integer}\NormalTok{(n),}
      \DataTypeTok{eval =} \KeywordTok{as.double}\NormalTok{ (a),}
      \DataTypeTok{evec =} \KeywordTok{as.double}\NormalTok{ (}\KeywordTok{rep}\NormalTok{(}\DecValTok{0}\NormalTok{, n }\OperatorTok{*}\StringTok{ }\NormalTok{n)),}
      \KeywordTok{as.double}\NormalTok{ (}\KeywordTok{rep}\NormalTok{(}\DecValTok{0}\NormalTok{, n)),}
      \KeywordTok{as.double}\NormalTok{ (}\KeywordTok{rep}\NormalTok{(}\DecValTok{0}\NormalTok{, n)),}
      \KeywordTok{as.integer}\NormalTok{(itmax),}
      \KeywordTok{as.double}\NormalTok{(eps)}
\NormalTok{    )}
  \KeywordTok{return}\NormalTok{ (}\KeywordTok{list}\NormalTok{ (}\DataTypeTok{eval =}\NormalTok{ h}\OperatorTok{$}\NormalTok{eval[r],}
                \DataTypeTok{evec =} \KeywordTok{matrix}\NormalTok{(h}\OperatorTok{$}\NormalTok{evec, n, n)))}
\NormalTok{\}}


\NormalTok{triangle2matrix <-}\StringTok{ }\ControlFlowTok{function}\NormalTok{ (x) \{}
\NormalTok{  m <-}\StringTok{ }\KeywordTok{length}\NormalTok{ (x)}
\NormalTok{  n <-}\StringTok{ }\KeywordTok{round}\NormalTok{ ((}\KeywordTok{sqrt}\NormalTok{ (}\DecValTok{1} \OperatorTok{+}\StringTok{ }\DecValTok{8} \OperatorTok{*}\StringTok{ }\NormalTok{m) }\OperatorTok{-}\StringTok{ }\DecValTok{1}\NormalTok{) }\OperatorTok{/}\StringTok{ }\DecValTok{2}\NormalTok{)}
\NormalTok{  h <-}
\StringTok{    }\KeywordTok{.C}\NormalTok{(}\StringTok{"trimat"}\NormalTok{, }\KeywordTok{as.integer}\NormalTok{ (n), }\KeywordTok{as.double}\NormalTok{ (x), }\KeywordTok{as.double}\NormalTok{ (}\KeywordTok{rep}\NormalTok{ (}\DecValTok{0}\NormalTok{, n }\OperatorTok{*}\StringTok{ }\NormalTok{n)))}
  \KeywordTok{return}\NormalTok{ (}\KeywordTok{matrix}\NormalTok{(h[[}\DecValTok{3}\NormalTok{]], n, n))}
\NormalTok{\}}

\NormalTok{matrix2triangle <-}\StringTok{ }\ControlFlowTok{function}\NormalTok{ (x) \{}
\NormalTok{  n <-}\StringTok{ }\KeywordTok{dim}\NormalTok{(x)[}\DecValTok{1}\NormalTok{]}
\NormalTok{  m <-}\StringTok{ }\NormalTok{n }\OperatorTok{*}\StringTok{ }\NormalTok{(n }\OperatorTok{+}\StringTok{ }\DecValTok{1}\NormalTok{) }\OperatorTok{/}\StringTok{ }\DecValTok{2}
\NormalTok{  h <-}
\StringTok{    }\KeywordTok{.C}\NormalTok{(}\StringTok{"mattri"}\NormalTok{, }\KeywordTok{as.integer}\NormalTok{ (n), }\KeywordTok{as.double}\NormalTok{ (x), }\KeywordTok{as.double}\NormalTok{ (}\KeywordTok{rep}\NormalTok{ (}\DecValTok{0}\NormalTok{, m)))}
  \KeywordTok{return}\NormalTok{ (h[[}\DecValTok{3}\NormalTok{]])}
\NormalTok{\}}

\NormalTok{trianglePrint <-}\StringTok{ }\ControlFlowTok{function}\NormalTok{ (x, }\DataTypeTok{width =} \DecValTok{6}\NormalTok{, }\DataTypeTok{precision =} \DecValTok{4}\NormalTok{) \{}
\NormalTok{  n <-}\StringTok{ }\KeywordTok{round}\NormalTok{ ((}\KeywordTok{sqrt}\NormalTok{ (}\DecValTok{1} \OperatorTok{+}\StringTok{ }\DecValTok{8} \OperatorTok{*}\StringTok{ }\NormalTok{m) }\OperatorTok{-}\StringTok{ }\DecValTok{1}\NormalTok{) }\OperatorTok{/}\StringTok{ }\DecValTok{2}\NormalTok{)}
\NormalTok{  h <-}\StringTok{ }\KeywordTok{.C}\NormalTok{(}\StringTok{"pritru"}\NormalTok{, }\KeywordTok{as.integer}\NormalTok{(n), }\KeywordTok{as.integer}\NormalTok{(width), }\KeywordTok{as.integer}\NormalTok{(precision), }\KeywordTok{as.double}\NormalTok{ (x))}
\NormalTok{\}}

\NormalTok{matrixPrint <-}\StringTok{ }\ControlFlowTok{function}\NormalTok{ (x, }\DataTypeTok{width =} \DecValTok{6}\NormalTok{, }\DataTypeTok{precision =} \DecValTok{4}\NormalTok{) \{}
\NormalTok{  n <-}\StringTok{ }\KeywordTok{nrow}\NormalTok{ (x)}
\NormalTok{  m <-}\StringTok{ }\KeywordTok{ncol}\NormalTok{ (x)}
\NormalTok{  h <-}\StringTok{ }\KeywordTok{.C}\NormalTok{(}\StringTok{"primat"}\NormalTok{, }\KeywordTok{as.integer}\NormalTok{(n), }\KeywordTok{as.integer}\NormalTok{(m), }\KeywordTok{as.integer}\NormalTok{(width), }\KeywordTok{as.integer}\NormalTok{(precision), }\KeywordTok{as.double}\NormalTok{ (x))}
\NormalTok{\}}

\NormalTok{jacobiR <-}
\StringTok{  }\ControlFlowTok{function}\NormalTok{(a,}
           \DataTypeTok{eps1 =} \FloatTok{1e-10}\NormalTok{,}
           \DataTypeTok{eps2 =} \FloatTok{1e-6}\NormalTok{,}
           \DataTypeTok{itmax =} \DecValTok{100}\NormalTok{,}
           \DataTypeTok{vectors =} \OtherTok{TRUE}\NormalTok{,}
           \DataTypeTok{verbose =} \OtherTok{FALSE}\NormalTok{) \{}
\NormalTok{    n <-}\StringTok{ }\KeywordTok{nrow}\NormalTok{(a)}
\NormalTok{    k <-}\StringTok{ }\KeywordTok{diag}\NormalTok{(n)}
\NormalTok{    itel <-}\StringTok{ }\DecValTok{1}
\NormalTok{    mx <-}\StringTok{ }\DecValTok{0}
\NormalTok{    saa <-}\StringTok{ }\KeywordTok{sum}\NormalTok{(a }\OperatorTok{^}\StringTok{ }\DecValTok{2}\NormalTok{)}
    \ControlFlowTok{repeat}\NormalTok{ \{}
      \ControlFlowTok{for}\NormalTok{ (i }\ControlFlowTok{in} \DecValTok{1}\OperatorTok{:}\NormalTok{(n }\OperatorTok{-}\StringTok{ }\DecValTok{1}\NormalTok{))}
        \ControlFlowTok{for}\NormalTok{ (j }\ControlFlowTok{in}\NormalTok{ (i }\OperatorTok{+}\StringTok{ }\DecValTok{1}\NormalTok{)}\OperatorTok{:}\NormalTok{n) \{}
\NormalTok{          aij <-}\StringTok{ }\NormalTok{a[i, j]}
\NormalTok{          bij <-}\StringTok{ }\KeywordTok{abs}\NormalTok{(aij)}
          \ControlFlowTok{if}\NormalTok{ (bij }\OperatorTok{<}\StringTok{ }\NormalTok{eps1)}
            \ControlFlowTok{next}\NormalTok{()}
\NormalTok{          mx <-}\StringTok{ }\KeywordTok{max}\NormalTok{(mx, bij)}
\NormalTok{          am <-}\StringTok{ }\NormalTok{(a[i, i] }\OperatorTok{-}\StringTok{ }\NormalTok{a[j, j]) }\OperatorTok{/}\StringTok{ }\DecValTok{2}
\NormalTok{          u <-}\StringTok{ }\KeywordTok{c}\NormalTok{(aij, }\OperatorTok{-}\NormalTok{am)}
\NormalTok{          u <-}\StringTok{ }\NormalTok{u }\OperatorTok{/}\StringTok{ }\KeywordTok{sqrt}\NormalTok{(}\KeywordTok{sum}\NormalTok{(u }\OperatorTok{^}\StringTok{ }\DecValTok{2}\NormalTok{))}
\NormalTok{          c <-}\StringTok{ }\KeywordTok{sqrt}\NormalTok{((}\DecValTok{1} \OperatorTok{+}\StringTok{ }\NormalTok{u[}\DecValTok{2}\NormalTok{]) }\OperatorTok{/}\StringTok{ }\DecValTok{2}\NormalTok{)}
\NormalTok{          s <-}\StringTok{ }\KeywordTok{sign}\NormalTok{(u[}\DecValTok{1}\NormalTok{]) }\OperatorTok{*}\StringTok{ }\KeywordTok{sqrt}\NormalTok{((}\DecValTok{1} \OperatorTok{-}\StringTok{ }\NormalTok{u[}\DecValTok{2}\NormalTok{]) }\OperatorTok{/}\StringTok{ }\DecValTok{2}\NormalTok{)}
\NormalTok{          ss <-}\StringTok{ }\NormalTok{s }\OperatorTok{^}\StringTok{ }\DecValTok{2}
\NormalTok{          cc <-}\StringTok{ }\NormalTok{c }\OperatorTok{^}\StringTok{ }\DecValTok{2}
\NormalTok{          sc <-}\StringTok{ }\NormalTok{s }\OperatorTok{*}\StringTok{ }\NormalTok{c}
\NormalTok{          ai <-}\StringTok{ }\NormalTok{a[i, ]}
\NormalTok{          aj <-}\StringTok{ }\NormalTok{a[j, ]}
\NormalTok{          aii <-}\StringTok{ }\NormalTok{a[i, i]}
\NormalTok{          ajj <-}\StringTok{ }\NormalTok{a[j, j]}
\NormalTok{          a[i, ] <-}\StringTok{ }\NormalTok{a[, i] <-}\StringTok{ }\NormalTok{c }\OperatorTok{*}\StringTok{ }\NormalTok{ai }\OperatorTok{-}\StringTok{ }\NormalTok{s }\OperatorTok{*}\StringTok{ }\NormalTok{aj}
\NormalTok{          a[j, ] <-}\StringTok{ }\NormalTok{a[, j] <-}\StringTok{ }\NormalTok{s }\OperatorTok{*}\StringTok{ }\NormalTok{ai }\OperatorTok{+}\StringTok{ }\NormalTok{c }\OperatorTok{*}\StringTok{ }\NormalTok{aj}
\NormalTok{          a[i, j] <-}\StringTok{ }\NormalTok{a[j, i] <-}\StringTok{ }\DecValTok{0}
\NormalTok{          a[i, i] <-}\StringTok{ }\NormalTok{aii }\OperatorTok{*}\StringTok{ }\NormalTok{cc }\OperatorTok{+}\StringTok{ }\NormalTok{ajj }\OperatorTok{*}\StringTok{ }\NormalTok{ss }\OperatorTok{-}\StringTok{ }\DecValTok{2} \OperatorTok{*}\StringTok{ }\NormalTok{sc }\OperatorTok{*}\StringTok{ }\NormalTok{aij}
\NormalTok{          a[j, j] <-}\StringTok{ }\NormalTok{ajj }\OperatorTok{*}\StringTok{ }\NormalTok{cc }\OperatorTok{+}\StringTok{ }\NormalTok{aii }\OperatorTok{*}\StringTok{ }\NormalTok{ss }\OperatorTok{+}\StringTok{ }\DecValTok{2} \OperatorTok{*}\StringTok{ }\NormalTok{sc }\OperatorTok{*}\StringTok{ }\NormalTok{aij}
          \ControlFlowTok{if}\NormalTok{ (vectors) \{}
\NormalTok{            ki <-}\StringTok{ }\NormalTok{k[, i]}
\NormalTok{            kj <-}\StringTok{ }\NormalTok{k[, j]}
\NormalTok{            k[, i] <-}\StringTok{ }\NormalTok{c }\OperatorTok{*}\StringTok{ }\NormalTok{ki }\OperatorTok{-}\StringTok{ }\NormalTok{s }\OperatorTok{*}\StringTok{ }\NormalTok{kj}
\NormalTok{            k[, j] <-}\StringTok{ }\NormalTok{s }\OperatorTok{*}\StringTok{ }\NormalTok{ki }\OperatorTok{+}\StringTok{ }\NormalTok{c }\OperatorTok{*}\StringTok{ }\NormalTok{kj}
\NormalTok{          \}}
\NormalTok{        \}}
\NormalTok{      ff <-}\StringTok{ }\KeywordTok{sqrt}\NormalTok{(saa }\OperatorTok{-}\StringTok{ }\KeywordTok{sum}\NormalTok{(}\KeywordTok{diag}\NormalTok{(a) }\OperatorTok{^}\StringTok{ }\DecValTok{2}\NormalTok{))}
      \ControlFlowTok{if}\NormalTok{ (verbose)}
        \KeywordTok{cat}\NormalTok{(}
          \StringTok{"Iteration "}\NormalTok{,}
          \KeywordTok{formatC}\NormalTok{(itel, }\DataTypeTok{digits =} \DecValTok{4}\NormalTok{),}
          \StringTok{"maxel "}\NormalTok{,}
          \KeywordTok{formatC}\NormalTok{(mx, }\DataTypeTok{width =} \DecValTok{10}\NormalTok{),}
          \StringTok{"loss "}\NormalTok{,}
          \KeywordTok{formatC}\NormalTok{(ff, }\DataTypeTok{width =} \DecValTok{10}\NormalTok{),}
          \StringTok{"}\CharTok{\textbackslash{}n}\StringTok{"}
\NormalTok{        )}
      \ControlFlowTok{if}\NormalTok{ ((mx }\OperatorTok{<}\StringTok{ }\NormalTok{eps1) }\OperatorTok{||}\StringTok{ }\NormalTok{(ff }\OperatorTok{<}\StringTok{ }\NormalTok{eps2) }\OperatorTok{||}\StringTok{ }\NormalTok{(itel }\OperatorTok{==}\StringTok{ }\NormalTok{itmax))}
        \ControlFlowTok{break}\NormalTok{()}
\NormalTok{      itel <-}\StringTok{ }\NormalTok{itel }\OperatorTok{+}\StringTok{ }\DecValTok{1}
\NormalTok{      mx <-}\StringTok{ }\DecValTok{0}
\NormalTok{    \}}
\NormalTok{    d <-}\StringTok{ }\KeywordTok{diag}\NormalTok{(a)}
\NormalTok{    o <-}\StringTok{ }\KeywordTok{order}\NormalTok{(d, }\DataTypeTok{decreasing =} \OtherTok{TRUE}\NormalTok{)}
    \ControlFlowTok{if}\NormalTok{ (vectors)}
      \KeywordTok{return}\NormalTok{(}\KeywordTok{list}\NormalTok{(}\DataTypeTok{values =}\NormalTok{ d[o], }\DataTypeTok{vectors =}\NormalTok{ k[, o]))}
    \ControlFlowTok{else}
      \KeywordTok{return}\NormalTok{(}\DataTypeTok{values =}\NormalTok{ d[o])}
\NormalTok{  \}}
\end{Highlighting}
\end{Shaded}

\subsection{C code}\label{c-code}

\subsubsection{jacobi.h}\label{jacobi.h}

\begin{Shaded}
\begin{Highlighting}[]
\PreprocessorTok{#ifndef JACOBI_H}
\PreprocessorTok{#define JACOBI_H}

\PreprocessorTok{#include }\ImportTok{<lapacke.h>}
\PreprocessorTok{#include }\ImportTok{<math.h>}
\PreprocessorTok{#include }\ImportTok{<stdbool.h>}
\PreprocessorTok{#include }\ImportTok{<stdio.h>}
\PreprocessorTok{#include }\ImportTok{<stdlib.h>}

\DataTypeTok{static} \KeywordTok{inline} \DataTypeTok{int}\NormalTok{ VINDEX(}\DataTypeTok{const} \DataTypeTok{int}\NormalTok{);}
\DataTypeTok{static} \KeywordTok{inline} \DataTypeTok{int}\NormalTok{ MINDEX(}\DataTypeTok{const} \DataTypeTok{int}\NormalTok{, }\DataTypeTok{const} \DataTypeTok{int}\NormalTok{, }\DataTypeTok{const} \DataTypeTok{int}\NormalTok{);}
\DataTypeTok{static} \KeywordTok{inline} \DataTypeTok{int}\NormalTok{ SINDEX(}\DataTypeTok{const} \DataTypeTok{int}\NormalTok{, }\DataTypeTok{const} \DataTypeTok{int}\NormalTok{, }\DataTypeTok{const} \DataTypeTok{int}\NormalTok{);}
\DataTypeTok{static} \KeywordTok{inline} \DataTypeTok{int}\NormalTok{ TINDEX(}\DataTypeTok{const} \DataTypeTok{int}\NormalTok{, }\DataTypeTok{const} \DataTypeTok{int}\NormalTok{, }\DataTypeTok{const} \DataTypeTok{int}\NormalTok{);}
\DataTypeTok{static} \KeywordTok{inline} \DataTypeTok{int}\NormalTok{ AINDEX(}\DataTypeTok{const} \DataTypeTok{int}\NormalTok{, }\DataTypeTok{const} \DataTypeTok{int}\NormalTok{, }\DataTypeTok{const} \DataTypeTok{int}\NormalTok{, }\DataTypeTok{const} \DataTypeTok{int}\NormalTok{, }\DataTypeTok{const} \DataTypeTok{int}\NormalTok{);}

\DataTypeTok{static} \KeywordTok{inline} \DataTypeTok{double}\NormalTok{ SQUARE(}\DataTypeTok{const} \DataTypeTok{double}\NormalTok{);}
\DataTypeTok{static} \KeywordTok{inline} \DataTypeTok{double}\NormalTok{ THIRD(}\DataTypeTok{const} \DataTypeTok{double}\NormalTok{);}
\DataTypeTok{static} \KeywordTok{inline} \DataTypeTok{double}\NormalTok{ FOURTH(}\DataTypeTok{const} \DataTypeTok{double}\NormalTok{);}
\DataTypeTok{static} \KeywordTok{inline} \DataTypeTok{double}\NormalTok{ FIFTH(}\DataTypeTok{const} \DataTypeTok{double}\NormalTok{);}

\DataTypeTok{static} \KeywordTok{inline} \DataTypeTok{double}\NormalTok{ MAX(}\DataTypeTok{const} \DataTypeTok{double}\NormalTok{, }\DataTypeTok{const} \DataTypeTok{double}\NormalTok{);}
\DataTypeTok{static} \KeywordTok{inline} \DataTypeTok{double}\NormalTok{ MIN(}\DataTypeTok{const} \DataTypeTok{double}\NormalTok{, }\DataTypeTok{const} \DataTypeTok{double}\NormalTok{);}
\DataTypeTok{static} \KeywordTok{inline} \DataTypeTok{int}\NormalTok{ IMIN(}\DataTypeTok{const} \DataTypeTok{int}\NormalTok{, }\DataTypeTok{const} \DataTypeTok{int}\NormalTok{);}
\DataTypeTok{static} \KeywordTok{inline} \DataTypeTok{int}\NormalTok{ IMAX(}\DataTypeTok{const} \DataTypeTok{int}\NormalTok{, }\DataTypeTok{const} \DataTypeTok{int}\NormalTok{);}

\DataTypeTok{static} \KeywordTok{inline} \DataTypeTok{int}\NormalTok{ VINDEX(}\DataTypeTok{const} \DataTypeTok{int}\NormalTok{ i) \{ }\ControlFlowTok{return}\NormalTok{ i - }\DecValTok{1}\NormalTok{; \}}

\DataTypeTok{static} \KeywordTok{inline} \DataTypeTok{int}\NormalTok{ MINDEX(}\DataTypeTok{const} \DataTypeTok{int}\NormalTok{ i, }\DataTypeTok{const} \DataTypeTok{int}\NormalTok{ j, }\DataTypeTok{const} \DataTypeTok{int}\NormalTok{ n) \{}
    \ControlFlowTok{return}\NormalTok{ (i - }\DecValTok{1}\NormalTok{) + (j - }\DecValTok{1}\NormalTok{) * n;}
\NormalTok{\}}

\DataTypeTok{static} \KeywordTok{inline} \DataTypeTok{int}\NormalTok{ AINDEX(}\DataTypeTok{const} \DataTypeTok{int}\NormalTok{ i, }\DataTypeTok{const} \DataTypeTok{int}\NormalTok{ j, }\DataTypeTok{const} \DataTypeTok{int}\NormalTok{ k, }\DataTypeTok{const} \DataTypeTok{int}\NormalTok{ n,}
                         \DataTypeTok{const} \DataTypeTok{int}\NormalTok{ m) \{}
    \ControlFlowTok{return}\NormalTok{ (i - }\DecValTok{1}\NormalTok{) + (j - }\DecValTok{1}\NormalTok{) * n + (k - }\DecValTok{1}\NormalTok{) * n * m;}
\NormalTok{\}}

\DataTypeTok{static} \KeywordTok{inline} \DataTypeTok{int}\NormalTok{ SINDEX(}\DataTypeTok{const} \DataTypeTok{int}\NormalTok{ i, }\DataTypeTok{const} \DataTypeTok{int}\NormalTok{ j, }\DataTypeTok{const} \DataTypeTok{int}\NormalTok{ n) \{}
    \ControlFlowTok{return}\NormalTok{ ((j - }\DecValTok{1}\NormalTok{) * n) - (j * (j - }\DecValTok{1}\NormalTok{) / }\DecValTok{2}\NormalTok{) + (i - j) - }\DecValTok{1}\NormalTok{;}
\NormalTok{\}}

\DataTypeTok{static} \KeywordTok{inline} \DataTypeTok{int}\NormalTok{ TINDEX(}\DataTypeTok{const} \DataTypeTok{int}\NormalTok{ i, }\DataTypeTok{const} \DataTypeTok{int}\NormalTok{ j, }\DataTypeTok{const} \DataTypeTok{int}\NormalTok{ n) \{}
    \ControlFlowTok{return}\NormalTok{ ((j - }\DecValTok{1}\NormalTok{) * n) - ((j - }\DecValTok{1}\NormalTok{) * (j - }\DecValTok{2}\NormalTok{) / }\DecValTok{2}\NormalTok{) + (i - (j - }\DecValTok{1}\NormalTok{)) - }\DecValTok{1}\NormalTok{;}
\NormalTok{\}}

\DataTypeTok{static} \KeywordTok{inline} \DataTypeTok{double}\NormalTok{ SQUARE(}\DataTypeTok{const} \DataTypeTok{double}\NormalTok{ x) \{ }\ControlFlowTok{return}\NormalTok{ x * x; \}}
\DataTypeTok{static} \KeywordTok{inline} \DataTypeTok{double}\NormalTok{ THIRD(}\DataTypeTok{const} \DataTypeTok{double}\NormalTok{ x) \{ }\ControlFlowTok{return}\NormalTok{ x * x * x; \}}
\DataTypeTok{static} \KeywordTok{inline} \DataTypeTok{double}\NormalTok{ FOURTH(}\DataTypeTok{const} \DataTypeTok{double}\NormalTok{ x) \{ }\ControlFlowTok{return}\NormalTok{ x * x * x * x; \}}
\DataTypeTok{static} \KeywordTok{inline} \DataTypeTok{double}\NormalTok{ FIFTH(}\DataTypeTok{const} \DataTypeTok{double}\NormalTok{ x) \{ }\ControlFlowTok{return}\NormalTok{ x * x * x * x * x; \}}

\DataTypeTok{static} \KeywordTok{inline} \DataTypeTok{double}\NormalTok{ MAX(}\DataTypeTok{const} \DataTypeTok{double}\NormalTok{ x, }\DataTypeTok{const} \DataTypeTok{double}\NormalTok{ y) \{}
    \ControlFlowTok{return}\NormalTok{ (x > y) ? x : y;}
\NormalTok{\}}

\DataTypeTok{static} \KeywordTok{inline} \DataTypeTok{double}\NormalTok{ MIN(}\DataTypeTok{const} \DataTypeTok{double}\NormalTok{ x, }\DataTypeTok{const} \DataTypeTok{double}\NormalTok{ y) \{}
    \ControlFlowTok{return}\NormalTok{ (x < y) ? x : y;}
\NormalTok{\}}

\DataTypeTok{static} \KeywordTok{inline} \DataTypeTok{int}\NormalTok{ IMAX(}\DataTypeTok{const} \DataTypeTok{int}\NormalTok{ x, }\DataTypeTok{const} \DataTypeTok{int}\NormalTok{ y) \{ }\ControlFlowTok{return}\NormalTok{ (x > y) ? x : y; \}}

\DataTypeTok{static} \KeywordTok{inline} \DataTypeTok{int}\NormalTok{ IMIN(}\DataTypeTok{const} \DataTypeTok{int}\NormalTok{ x, }\DataTypeTok{const} \DataTypeTok{int}\NormalTok{ y) \{ }\ControlFlowTok{return}\NormalTok{ (x < y) ? x : y; \}}

\PreprocessorTok{#endif }\CommentTok{/* JACOBI_H */}
\end{Highlighting}
\end{Shaded}

\subsubsection{jacobi.c}\label{jacobi.c}

\begin{Shaded}
\begin{Highlighting}[]
\PreprocessorTok{#include }\ImportTok{"jacobi.h"}

\DataTypeTok{void}\NormalTok{ jacobiC(}\DataTypeTok{const} \DataTypeTok{int}\NormalTok{ *nn, }\DataTypeTok{double}\NormalTok{ *a, }\DataTypeTok{double}\NormalTok{ *evec, }\DataTypeTok{double}\NormalTok{ *oldi, }\DataTypeTok{double}\NormalTok{ *oldj,}
             \DataTypeTok{int}\NormalTok{ *itmax, }\DataTypeTok{double}\NormalTok{ *eps) \{}
    \DataTypeTok{int}\NormalTok{ n = *nn, itel = }\DecValTok{1}\NormalTok{;}
    \DataTypeTok{double}\NormalTok{ d = }\FloatTok{0.0}\NormalTok{, s = }\FloatTok{0.0}\NormalTok{, t = }\FloatTok{0.0}\NormalTok{, u = }\FloatTok{0.0}\NormalTok{, v = }\FloatTok{0.0}\NormalTok{, p = }\FloatTok{0.0}\NormalTok{, q = }\FloatTok{0.0}\NormalTok{,}
\NormalTok{           r = }\FloatTok{0.0}\NormalTok{;}
    \DataTypeTok{double}\NormalTok{ fold = }\FloatTok{0.0}\NormalTok{, fnew = }\FloatTok{0.0}\NormalTok{;}
    \ControlFlowTok{for}\NormalTok{ (}\DataTypeTok{int}\NormalTok{ i = }\DecValTok{1}\NormalTok{; i <= n; i++) \{}
        \ControlFlowTok{for}\NormalTok{ (}\DataTypeTok{int}\NormalTok{ j = }\DecValTok{1}\NormalTok{; j <= n; j++) \{}
\NormalTok{            evec[MINDEX(i, j, n)] = (i == j) ? }\FloatTok{1.0}\NormalTok{ : }\FloatTok{0.0}\NormalTok{;}
\NormalTok{        \}}
\NormalTok{    \}}
    \ControlFlowTok{for}\NormalTok{ (}\DataTypeTok{int}\NormalTok{ i = }\DecValTok{1}\NormalTok{; i <= n; i++) \{}
\NormalTok{        fold += SQUARE(a[TINDEX(i, i, n)]);}
\NormalTok{    \}}
    \ControlFlowTok{while}\NormalTok{ (true) \{}
        \ControlFlowTok{for}\NormalTok{ (}\DataTypeTok{int}\NormalTok{ j = }\DecValTok{1}\NormalTok{; j <= n - }\DecValTok{1}\NormalTok{; j++) \{}
            \ControlFlowTok{for}\NormalTok{ (}\DataTypeTok{int}\NormalTok{ i = j + }\DecValTok{1}\NormalTok{; i <= n; i++) \{}
\NormalTok{                p = a[TINDEX(i, j, n)];}
\NormalTok{                q = a[TINDEX(i, i, n)];}
\NormalTok{                r = a[TINDEX(j, j, n)];}
                \ControlFlowTok{if}\NormalTok{ (fabs(p) < }\FloatTok{1e-10}\NormalTok{) }\ControlFlowTok{continue}\NormalTok{;}
\NormalTok{                d = (q - r) / }\FloatTok{2.0}\NormalTok{;}
\NormalTok{                s = (p < }\DecValTok{0}\NormalTok{) ? -}\FloatTok{1.0}\NormalTok{ : }\FloatTok{1.0}\NormalTok{;}
\NormalTok{                t = -d / sqrt(SQUARE(d) + SQUARE(p));}
\NormalTok{                u = sqrt((}\DecValTok{1}\NormalTok{ + t) / }\DecValTok{2}\NormalTok{);}
\NormalTok{                v = s * sqrt((}\DecValTok{1}\NormalTok{ - t) / }\DecValTok{2}\NormalTok{);}
                \ControlFlowTok{for}\NormalTok{ (}\DataTypeTok{int}\NormalTok{ k = }\DecValTok{1}\NormalTok{; k <= n; k++) \{}
                    \DataTypeTok{int}\NormalTok{ ik = IMIN(i, k);}
                    \DataTypeTok{int}\NormalTok{ ki = IMAX(i, k);}
                    \DataTypeTok{int}\NormalTok{ jk = IMIN(j, k);}
                    \DataTypeTok{int}\NormalTok{ kj = IMAX(j, k);}
\NormalTok{                    oldi[VINDEX(k)] = a[TINDEX(ki, ik, n)];}
\NormalTok{                    oldj[VINDEX(k)] = a[TINDEX(kj, jk, n)];}
\NormalTok{                \}}
                \ControlFlowTok{for}\NormalTok{ (}\DataTypeTok{int}\NormalTok{ k = }\DecValTok{1}\NormalTok{; k <= n; k++) \{}
                    \DataTypeTok{int}\NormalTok{ ik = IMIN(i, k);}
                    \DataTypeTok{int}\NormalTok{ ki = IMAX(i, k);}
                    \DataTypeTok{int}\NormalTok{ jk = IMIN(j, k);}
                    \DataTypeTok{int}\NormalTok{ kj = IMAX(j, k);}
\NormalTok{                    a[TINDEX(ki, ik, n)] =}
\NormalTok{                        u * oldi[VINDEX(k)] - v * oldj[VINDEX(k)];}
\NormalTok{                    a[TINDEX(kj, jk, n)] =}
\NormalTok{                        v * oldi[VINDEX(k)] + u * oldj[VINDEX(k)];}
\NormalTok{                \}}
                \ControlFlowTok{for}\NormalTok{ (}\DataTypeTok{int}\NormalTok{ k = }\DecValTok{1}\NormalTok{; k <= n; k++) \{}
\NormalTok{                    oldi[VINDEX(k)] = evec[MINDEX(k, i, n)];}
\NormalTok{                    oldj[VINDEX(k)] = evec[MINDEX(k, j, n)];}
\NormalTok{                    evec[MINDEX(k, i, n)] =}
\NormalTok{                        u * oldi[VINDEX(k)] - v * oldj[VINDEX(k)];}
\NormalTok{                    evec[MINDEX(k, j, n)] =}
\NormalTok{                        v * oldi[VINDEX(k)] + u * oldj[VINDEX(k)];}
\NormalTok{                \}}
\NormalTok{                a[TINDEX(i, i, n)] =}
\NormalTok{                    SQUARE(u) * q + SQUARE(v) * r - }\DecValTok{2}\NormalTok{ * u * v * p;}
\NormalTok{                a[TINDEX(j, j, n)] =}
\NormalTok{                    SQUARE(v) * q + SQUARE(u) * r + }\DecValTok{2}\NormalTok{ * u * v * p;}
\NormalTok{                a[TINDEX(i, j, n)] =}
\NormalTok{                    u * v * (q - r) + (SQUARE(u) - SQUARE(v)) * p;}
\NormalTok{            \}}
\NormalTok{        \}}
\NormalTok{        fnew = }\FloatTok{0.0}\NormalTok{;}
        \ControlFlowTok{for}\NormalTok{ (}\DataTypeTok{int}\NormalTok{ i = }\DecValTok{1}\NormalTok{; i <= n; i++) \{}
\NormalTok{            fnew += SQUARE(a[TINDEX(i, i, n)]);}
\NormalTok{        \}}
        \ControlFlowTok{if}\NormalTok{ (((fnew - fold) < *eps) || (itel == *itmax)) }\ControlFlowTok{break}\NormalTok{;}
\NormalTok{        fold = fnew;}
\NormalTok{        itel++;}
\NormalTok{    \}}
    \ControlFlowTok{return}\NormalTok{;}
\NormalTok{\}}

\DataTypeTok{void}\NormalTok{ primat(}\DataTypeTok{const} \DataTypeTok{int}\NormalTok{ *n, }\DataTypeTok{const} \DataTypeTok{int}\NormalTok{ *m, }\DataTypeTok{const} \DataTypeTok{int}\NormalTok{ *w, }\DataTypeTok{const} \DataTypeTok{int}\NormalTok{ *p,}
            \DataTypeTok{const} \DataTypeTok{double}\NormalTok{ *x) \{}
    \ControlFlowTok{for}\NormalTok{ (}\DataTypeTok{int}\NormalTok{ i = }\DecValTok{1}\NormalTok{; i <= *n; i++) \{}
        \ControlFlowTok{for}\NormalTok{ (}\DataTypeTok{int}\NormalTok{ j = }\DecValTok{1}\NormalTok{; j <= *m; j++) \{}
\NormalTok{            printf(}\StringTok{" %*.*f "}\NormalTok{, *w, *p, x[MINDEX(i, j, *n)]);}
\NormalTok{        \}}
\NormalTok{        printf(}\StringTok{"}\SpecialCharTok{\textbackslash{}n}\StringTok{"}\NormalTok{);}
\NormalTok{    \}}
\NormalTok{    printf(}\StringTok{"}\SpecialCharTok{\textbackslash{}n\textbackslash{}n}\StringTok{"}\NormalTok{);}
    \ControlFlowTok{return}\NormalTok{;}
\NormalTok{\}}

\DataTypeTok{void}\NormalTok{ pritru(}\DataTypeTok{const} \DataTypeTok{int}\NormalTok{ *n, }\DataTypeTok{const} \DataTypeTok{int}\NormalTok{ *w, }\DataTypeTok{const} \DataTypeTok{int}\NormalTok{ *p, }\DataTypeTok{const} \DataTypeTok{double}\NormalTok{ *x) \{}
    \ControlFlowTok{for}\NormalTok{ (}\DataTypeTok{int}\NormalTok{ i = }\DecValTok{1}\NormalTok{; i <= *n; i++) \{}
        \ControlFlowTok{for}\NormalTok{ (}\DataTypeTok{int}\NormalTok{ j = }\DecValTok{1}\NormalTok{; j <= i; j++) \{}
\NormalTok{            printf(}\StringTok{" %*.*f "}\NormalTok{, *w, *p, x[TINDEX(i, j, *n)]);}
\NormalTok{        \}}
\NormalTok{        printf(}\StringTok{"}\SpecialCharTok{\textbackslash{}n}\StringTok{"}\NormalTok{);}
\NormalTok{    \}}
\NormalTok{    printf(}\StringTok{"}\SpecialCharTok{\textbackslash{}n\textbackslash{}n}\StringTok{"}\NormalTok{);}
    \ControlFlowTok{return}\NormalTok{;}
\NormalTok{\}}

\DataTypeTok{void}\NormalTok{ trimat(}\DataTypeTok{const} \DataTypeTok{int}\NormalTok{ *n, }\DataTypeTok{const} \DataTypeTok{double}\NormalTok{ *x, }\DataTypeTok{double}\NormalTok{ *y) \{}
    \DataTypeTok{int}\NormalTok{ nn = *n;}
    \ControlFlowTok{for}\NormalTok{ (}\DataTypeTok{int}\NormalTok{ i = }\DecValTok{1}\NormalTok{; i <= nn; i++) \{}
        \ControlFlowTok{for}\NormalTok{ (}\DataTypeTok{int}\NormalTok{ j = }\DecValTok{1}\NormalTok{; j <= nn; j++) \{}
\NormalTok{            y[MINDEX(i, j, nn)] =}
\NormalTok{                (i >= j) ? x[TINDEX(i, j, nn)] : x[TINDEX(j, i, nn)];}
\NormalTok{        \}}
\NormalTok{    \}}
    \ControlFlowTok{return}\NormalTok{;}
\NormalTok{\}}

\DataTypeTok{void}\NormalTok{ mattri(}\DataTypeTok{const} \DataTypeTok{int}\NormalTok{ *n, }\DataTypeTok{const} \DataTypeTok{double}\NormalTok{ *x, }\DataTypeTok{double}\NormalTok{ *y) \{}
    \DataTypeTok{int}\NormalTok{ nn = *n;}
    \ControlFlowTok{for}\NormalTok{ (}\DataTypeTok{int}\NormalTok{ j = }\DecValTok{1}\NormalTok{; j <= nn; j++) \{}
        \ControlFlowTok{for}\NormalTok{ (}\DataTypeTok{int}\NormalTok{ i = j; i <= nn; i++) \{}
\NormalTok{            y[TINDEX(i, j, nn)] = x[MINDEX(i, j, nn)];}
\NormalTok{        \}}
\NormalTok{    \}}
    \ControlFlowTok{return}\NormalTok{;}
\NormalTok{\}}
\end{Highlighting}
\end{Shaded}

\section*{References}\label{references}
\addcontentsline{toc}{section}{References}

\hypertarget{refs}{}
\hypertarget{ref-deleeuw_ferrari_R_08a}{}
De Leeuw, J., and D.B. Ferrari. 2008. ``Using Jacobi Plane Rotations in
R.'' Preprint Series 556. Los Angeles, CA: UCLA Department of
Statistics.
\url{http://deleeuwpdx.net/janspubs/2008/reports/deleeuw_R_08a.pdf}.


\end{document}
